\chapter*{Abstract}
\addcontentsline{toc}{chapter}{Abstract}

This is a broad survey of how the world's esoteric and spiritual traditions have modeled reality, the
universe, and the inner experience of being here. It is not a history and not an argument for any one faith.
It compares many models of the same subject.

For thousands of years, in many languages and on every continent, people devoted their lives to the same few
questions. What is the source of everything. How did the many come from the one. What is a human being made
of, on the inside. Why are we here, and where are we going. Each tradition answered in its own images and its
own technical language. From a distance the answers look very different. Read closely and laid side by side,
they share a common structure.

Two traditions are treated in the most depth, because they went the furthest in building a complete and
explicit system, the Jewish \textbf{Kabbalah} and \textbf{Tibetan Buddhism}. Around them sit six living
esoteric traditions, six ancient cosmologies, and eighteen indigenous and less-documented visions. Every
tradition is introduced the same way. First with its creation story told in plain modern words with no
special terms, so the bare shape of the model is visible. Then with its own concepts, structures, and
vocabulary.

The aim is practical as well as contemplative. To model consciousness, the inner life, the sense of the
divine, and the worlds of beings both bright and dark, it helps to know the full depth and scope of what the
traditions already saw. This survey gathers that scope in one place. It asks at the end whether all these maps
trace one territory, a single movement of the one source going out into the many and finding its way back to
itself.
